% Source: https://artofproblemsolving.com/community/c3h1444703_is_this_trivial
% Stable identity: aops-topic-1444703-q01
\textbf{(a)} \emph{In $x$.} Here $D=x^2-y^2$ has degree $2$ in $x$, and the remainder must have degree less than $2$ in $x$. Since $P=2(x^2-y^2)+y^2$ and $y^2$ has degree $0$ in $x$, the quotient is $2$ and the remainder is $\boxed{y^2}$.

\emph{In $y$.} Now $D=-y^2+x^2$ has degree $2$ in $y$, and the remainder must have degree less than $2$ in $y$. Since $P=-y^2+2x^2=1\cdot(-y^2+x^2)+x^2$ and $x^2$ has degree $0$ in $y$, the quotient is $1$ and the remainder is $\boxed{x^2}$.

\textbf{(b)} For polynomials in one variable, division by $D$ gives a unique quotient and remainder, the remainder being the part of degree less than $\deg D$. With two variables we must choose which variable's degree is used, and the two choices give different (but each correct and unique) answers: $y^2$ for division in $x$, and $x^2$ for division in $y$. The two answers are consistent, since they differ by a multiple of $D$:
\[
x^2-y^2=1\cdot D .
\]
Both of the student's identities are true; the first one is the division in $x$ and the second the division in $y$.
